\exercisesection{}

在下面的习题中，当 E 是巴拿赫空间时，我们用 \(\mathcal{C}alk(E)\) 表示卡尔金代数 \(\mathcal{L}(E)/\mathcal{L}^{c}(E)\) 。用 \(\pi\) 表示从 \(\mathcal{L}(E)\) 到 \(\mathcal{C}alk(E)\) 上的典范投影 。

若 E 是复希尔伯特空间，我们给 \(\mathcal{C}alk(\mathrm{E})\) 赋以由 \(\mathcal{L}(\mathrm{E})\) 的对合经取商而诱导的对合。

\begin{enumerate}
\def\labelenumi{\arabic{enumi})}
\tightlist
\item
  设 E 是局部凸空间。
\end{enumerate}

\begin{enumerate}
\def\labelenumi{\alph{enumi})}
\item
  假设 E 中 0 的每个邻域都包含 E 的一个有限余维数的向量子空间。证明：E 的弗雷德霍姆自同态所成的空间 \(\mathcal{F}(\mathrm{E})\) 在赋以有界收敛拓扑的 \(\mathcal{L}(\mathrm{E})\) 中不是闭的。
\item
  给出一个无穷维弗雷歇空间的例子，使得其中 0 的每个邻域都包含 E 的一个有限余维数的向量子空间。c) 设 \(E_{\sigma}\) 是赋以弱拓扑的无穷维巴拿赫空间。集合 \(\mathcal{F}(E_{\sigma})\) 在 \(\mathcal{L}(E_{\sigma})\) 中不是开的。
\end{enumerate}

\begin{enumerate}
\def\labelenumi{\arabic{enumi})}
\setcounter{enumi}{1}
\tightlist
\item
  设 E 是无穷维巴拿赫空间，且 \(u \in \mathcal{L}(\mathrm{E})\) 。我们用 \(\mathrm{S}(u)\) 表示如下的复数 \(\lambda\) 所成的集：存在实数 \(\delta > 0\) 和 E 的一个有限余维数子空间 F，使得 \(\| u(x) - \lambda x \| \geqslant \delta \| x \|\) 对所有 \(x \in \mathrm{F}\) 成立。
\end{enumerate}

\begin{enumerate}
\def\labelenumi{\alph{enumi})}
\item
  集合 \(S(u)\) 在 C 中是开的，且包含 u 的谱在 C 中的补集。
\item
  若 \(\lambda \in \mathrm{S}(u)\) ，则自同态 \(u - \lambda 1_{\mathrm{E}}\) 具有闭的像和有限维的核。
\item
  补集 \(\mathrm{F}(u) = \mathbf{C} - \mathrm{S}(u)\) 非空。（证明 \(\mathrm{Sp}_{\mathrm{e}}(u)\) 中一个模最大的元素属于 \(\mathrm{F}(u)\) 。）
\end{enumerate}

\begin{enumerate}
\def\labelenumi{\arabic{enumi})}
\setcounter{enumi}{2}
\tightlist
\item
  设 E 是无穷维复希尔伯特空间。
\end{enumerate}

\begin{enumerate}
\def\labelenumi{\alph{enumi})}
\item
  对合代数 \(\mathcal{C}alk(E)\) 是星代数。
\item
  设 u 是 \(\mathcal{L}(\mathrm{E})\) 的使 \(u^{*}u\) 为紧算子的元素。证明 u 是紧算子。
\item
  设 G 是 \(\mathcal{C}alk(\mathrm{E})\) 中可逆元素所成的群，\(G_{0}\) 是它的单位分量。设 \(p: G \to G/G_{0}\) 是典范投影。存在唯一的群同构 \(\alpha: Z \to G/G_{0}\) ，使得对 E 的每个弗雷德霍姆自同态 u 都有 \(\alpha(\operatorname{ind}(u)) = p(\pi(u))\) 。
\end{enumerate}

\begin{enumerate}
\def\labelenumi{\arabic{enumi})}
\setcounter{enumi}{3}
\tightlist
\item
  设 E 是复希尔伯特空间。设 \(\xi\) 是 \(\mathcal{C}alk(E)\) 的一个可逆元素，我们称 \(\xi\) 的指标为 E 的任一弗雷德霍姆自同态 \(u\) （满足 \(\pi(u) = \xi\) ）的指标。
\end{enumerate}

\begin{enumerate}
\def\labelenumi{\alph{enumi})}
\item
  证明：对每个埃尔米特的 \(\xi \in \mathcal{C}alk(\mathrm{E})\) ，存在埃尔米特的 \(u \in \mathcal{L}(\mathrm{E})\) 使得 \(\pi(u) = \xi\) 。
\item
  证明：对每个酉的且指标为零的 \(\xi \in \mathcal{C}alk(\mathrm{E})\) ，存在酉的 \(u \in \mathcal{L}(\mathrm{E})\) 使得 \(\pi(u) = \xi\) 。（设 \(v\) 可逆且满足 \(\pi(v) = \xi\) ；考察 \(v = j|v|\) 这一 \(v\) 的极分解，验证 \(|v| - 1_{\mathrm{E}}\) 是紧的，再推出 \(v - j\) 是紧的。）
\item
  证明：对每个酉的 \(\xi \in \mathcal{C}alk(\mathrm{E})\) （满足 \(\xi^2 = 1\) ），存在酉的 \(u \in \mathcal{L}(\mathrm{E})\) 使得 \(u^2 = 1\) 且 \(\pi(u) = \xi\) 。
\end{enumerate}

\(d)\) 对每个元素 \(p \in \mathcal{C}alk(\mathrm{E})\) ，只要 \(p^2 = p\) ，就存在投影 \(\mathrm{P} \in \mathcal{L}(\mathrm{E})\) 使得 \(\pi(\mathrm{P}) = p\) 。

\begin{enumerate}
\def\labelenumi{\arabic{enumi})}
\setcounter{enumi}{4}
\tightlist
\item
  设 E 是复希尔伯特空间。
\end{enumerate}

\begin{enumerate}
\def\labelenumi{\alph{enumi})}
\tightlist
\item
  设 \(n\) 是满足 \(n \geqslant 2\) 的整数。我们置 \(\mathrm{I} = \mathbf{N} \times \{0, \ldots, n\}\) 。设 \((\sigma_i)_{1 \leqslant i \leqslant n}\) 是 I 的由下式定义的置换族
\end{enumerate}

\[
\begin{array}{c} \sigma_ {i} (m, i) = (m - 1, i) \text {for} m \geqslant 1 \\ \sigma_ {i} (0, i) = (0, i - 1) \\ \sigma_ {i} (m, i - 1) = (m + 1, i - 1) \text {for} m \geqslant 0 \\ \sigma_ {i} (m, j) = (m, j) \text {for} m \geqslant 0 \text {and} j \notin \{i, i - 1 \}. \end{array}
\]

\begin{enumerate}
\def\labelenumi{\alph{enumi})}
\setcounter{enumi}{1}
\item
  用 \((e_{m,i})\) 表示 \(\ell^{2}(\mathrm{I})\) 的典范标准正交基。设 \(u_{i}\) 是 \(\ell^{2}(\mathrm{I})\) 的自同态，使得 \(u_{i}(e_{m,j}) = e_{\sigma_{i}(m,j)}\) 对所有 \((m,j) \in \mathrm{I}\) 成立。自同态 \(u_{i}\) 是酉的。
\item
  元素 \(\pi(u_i)\) 两两交换。
\end{enumerate}

\(d)\) 不存在 E 的由两两交换的酉自同态组成的族 \((v_{1},\ldots,v_{n})\) ，使得 \(\pi(v_{i})=\pi(u_{i})\) 对所有 \(i\) 成立。

\begin{enumerate}
\def\labelenumi{\arabic{enumi})}
\setcounter{enumi}{5}
\tightlist
\item
  设 E 是复希尔伯特空间。回顾 \(\mathcal{C}alk(E)\) 是星代数（习题 3）。
\end{enumerate}

  设 \(\mathfrak{F}\) 是 \(\mathbf{N}\) 上的一个非主超滤子。设 F 是 E 中沿 \(\mathfrak{F}\) 弱收敛于 0 的有界序列所成的向量空间，\(F_0\) 是 F 中由沿 \(\mathfrak{F}\) 依范数收敛于 0 的序列组成的子空间。我们置 \(\widetilde{\mathrm{H}} = \mathrm{F} / \mathrm{F}_0\) 。a) 对 F 的所有元素 \(x = (x_n)\) 和 \(y = (y_n)\) ，证明序列 \((\langle x_n | y_n \rangle)_{n \in \mathbf{N}}\) 沿 \(\mathfrak{F}\) 收敛。用 \(\langle x | y \rangle_{\mathrm{F}}\) 表示其极限。

\begin{enumerate}
\def\labelenumi{\alph{enumi})}
\setcounter{enumi}{1}
\item
  证明 \((x,y)\mapsto\langle x\mid y\rangle_{\mathrm{F}}\) 是 F 上的正埃尔米特形式，且 \(\langle x\mid x\rangle_{F}=0\) 当且仅当 \(x\in F_{0}\) ；因而这个埃尔米特形式在 \(\widetilde{H}\) 上定义了一个分离的预希尔伯特空间结构。
\item
  设 H 是由分离的预希尔伯特空间 \(\widetilde{H}\) 完备化得到的希尔伯特空间。证明：对 \(\mathcal{L}(E)\) 中的每个 u，都得到一个元素 \(\varrho(u)\) （属于 \(\mathcal{L}(H)\) ），这只需把 F 到 F 的线性映射 \((x_{n})\mapsto(u(x_{n}))\) 过渡到商。d) 映射 \(u\mapsto\varrho(u)\) 诱导 C*-代数的一个等距态射 \(\mathcal{C}alk(E)\to\mathcal{L}(H)\) 。
\end{enumerate}

（我们以后将看到，每个 C*-代数 A 都有一个等距表示 \(A \to \mathcal{L}(E)\) ，其中 E 是某个希尔伯特空间，参见 V, p. 442 的定理 2。）

\begin{enumerate}
\def\labelenumi{\arabic{enumi})}
\setcounter{enumi}{6}
\tightlist
\item
  我们给单位圆 \(S_{1} \subset \mathbf{C}\) 赋以规范化哈尔测度。设 E 是复希尔伯特空间 \(\mathrm{L}^{2}(\mathbf{S}_{1})\) 。对 \(n \in \mathbf{Z}\) ，设 \(e_{n}\) 是 E 的元素 \(z \mapsto z^{n}\) 。设 F ⊂ E 是由诸 \(e_{n}\) （其中 \(n \geqslant 0\) ）生成的闭子空间，并设 \(p: E \to F\) 是正交投影。对 \(f \in L^{\infty}(\mathbf{S}_{1})\) ，我们在 E（相应地，F）上定义线性映射 \(M_{f}\)（相应地，\(T_{f}\)）如下：\(\mathrm{M}_{f}(g) = fg\) （\(g \in E\) ），且 \(T_{f} = p \circ M_{f}\) 。我们称 \(T_{f}\) 为以 f 为符号的特普利茨算子。(a) 线性映射 \(M_{f}\) 与 \(T_{f}\) 连续。\(M_{f}\) 的谱是 f 的傅里叶系数 \(c_{n} = \langle f | e_{n} \rangle\) 所成之集的闭包。自同态 \(M_{f}\) 是弗雷德霍姆自同态，当且仅当它是里斯自同态。
\end{enumerate}

\begin{enumerate}
\def\labelenumi{\alph{enumi})}
\setcounter{enumi}{1}
\item
  自同态 \(T_{e_{n}}\) （其中 \(n \in Z\) ）都是弗雷德霍姆自同态，且 \(\text{ind}(\mathrm{T}_{e_{n}}) = -n\) 对所有 \(n \in Z\) 成立。
\item
  \(\mathrm{T}_{\overline{f}} = \mathrm{T}_f^*\) 对每个函数 \(f\in \mathrm{L}^{\infty}(\mathbf{S}_1)\) 都成立。
\item
  若 \(T_{f}\) 是同构，则 \(M_{f}\) 是同构，且 f 在 \(\mathrm{L}^{\infty}(\mathbf{S}_{1})\) 中可逆。由此推出 \(M_{f}\) 的谱包含于 \(T_{f}\) 的谱，且 \(\|T_{f}\|=\|f\|\) 。
\item
  设 \(f \in \mathcal{C}(\mathbf{S}_1)\) 是连续函数，\(g \in \mathrm{L}^{\infty}(\mathbf{S}_1)\) 。自同态 \(\mathrm{T}_f\mathrm{T}_g - \mathrm{T}_{fg}\) 与 \(\mathrm{T}_g\mathrm{T}_f - \mathrm{T}_{gf}\) 是紧算子。（先考虑 \(f = e_n\) 的情形。）
\item
  设 \(\mathcal{C}alk(\mathrm{F})\) 是 F 的卡尔金代数，\(\pi\colon\mathcal{L}(\mathrm{F})\to\mathcal{C}alk(\mathrm{F})\) 是典范投影。证明 \(f\mapsto\pi(\mathrm{T}_{f})\) 是从 \(\mathcal{C}(\mathbf{S}_{1})\) 到 \(\mathcal{C}alk(\mathrm{F})\) 中的对合代数的单态射。（为证单性，可利用 I, p. 31 的命题 1。）
\end{enumerate}

\(g)\) 设 \(f \in \mathcal{C}(\mathbf{S}_1)\) 。则 \(\mathrm{T}_f\) 是弗雷德霍姆自同态，当且仅当 \(f\) 在 \(\mathbf{S}_1\) 上不取零值。（利用 I, p. 106 的命题 5。）

\(h)\) 设 \(f \in \mathcal{C}(\mathbf{S}_1)\) 且 \(f\) 在 \(\mathbf{S}_1\) 上不取零值。\(\mathrm{T}_f\) 的指标等于 \(n \in \mathbf{Z}\) ，当且仅当道路 \(t \mapsto f(e^{2i\pi t})\) （在 \(\mathbf{C} - \{0\}\) 中）同伦于道路 \(t \mapsto e_n(t)\) （参见 TA, III, p. 229, 定义 1）。

\begin{enumerate}
\def\labelenumi{\arabic{enumi})}
\setcounter{enumi}{7}
\tightlist
\item
  我们沿用上一习题的记号。
\end{enumerate}

\begin{enumerate}
\def\labelenumi{\alph{enumi})}
\item
  用 \(\mathrm{T}(\mathbf{S}_{1})\) 表示 \(\mathcal{L}(\mathrm{F})\) 中由特普利茨算子 \(T_{e_{1}}\) 与 \(T_{e_{-1}}\) 生成的闭子代数。证明 \(\mathrm{T}(\mathbf{S}_{1})\) 是 \(\mathcal{L}(\mathrm{F})\) 的星子代数。
\item
  设 \(f \in \mathcal{C}(\mathbf{S}_1)\) 。证明 \(T_f\) 属于 \(T(\mathbf{S}_1)\) 。
\item
  \(\mathrm{T}(\mathbf{S}_{1})\) 在 F 上的表示是不可约的；星代数 \(\mathrm{T}(\mathbf{S}_{1})\) 包含 \(\mathcal{L}^{\mathrm{c}}(\mathrm{F})\) 。
\end{enumerate}

\(d)\) \(u \in \mathrm{T}(\mathbf{S}_1)\) 当且仅当存在 \(f \in \mathcal{C}(\mathbf{S}_1)\) 与 \(h \in \mathcal{L}^c(\mathrm{F})\) 使得 \(u = \mathrm{T}_f + h\) 。（证明形如 \(\mathrm{T}_f + h\) 的自同态所成之集在 \(\mathcal{L}(\mathrm{F})\) 中构成一个闭星代数。）

\begin{enumerate}
\def\labelenumi{\alph{enumi})}
\setcounter{enumi}{4}
\tightlist
\item
  存在对合代数态射的一个正合列
\end{enumerate}

\[
0 \longrightarrow \mathcal {L} ^ {\mathrm{c}} (\mathrm{F}) \longrightarrow \mathrm{T} (\mathbf {S} _ {1}) \stackrel {\varrho} {\longrightarrow} \mathcal {C} (\mathbf {S} _ {1}) \longrightarrow 0
\]

使得 \(\varrho(\mathrm{T}_f) = f\) 对每个函数 \(f \in \mathcal{C}(\mathbf{S}_1)\) 成立。

\begin{enumerate}
\def\labelenumi{\arabic{enumi})}
\setcounter{enumi}{8}
\item
  设 B 是无穷维巴拿赫空间。用 E 表示 B 的强对偶，用 F 表示赋以拓扑 \(\sigma(\mathrm{B}^{\prime},\mathrm{B})\) 的 B 的对偶。证明：存在从 E 到 F 中的连续双射线性映射 u，以及从 E 到 F 中的紧线性映射 h，使得 \(u + h\) 有无穷维的余核。
\item
  设 E 是分离的局部凸空间，u 与 h 是 \(\mathcal{L}(\mathrm{E})\) 中交换的元素。假设存在整数 n 使得 \(h^{n}\) 是紧自同态。
\end{enumerate}

\begin{enumerate}
\def\labelenumi{\alph{enumi})}
\item
  若 u 是弗雷德霍姆自同态，则 u-h 是弗雷德霍姆自同态。
\item
  若 \(u\) 是里斯自同态，则 \(u - h\) 是里斯自同态。
\end{enumerate}

\begin{enumerate}
\def\labelenumi{\arabic{enumi})}
\setcounter{enumi}{10}
\tightlist
\item
  设 \(\mathrm{E} = \ell^2 (\mathbf{N})\) ，\((e_n)\) 是 E 的典范基。我们定义 \(u\in \mathcal{L}(\mathrm{E})\) 为满足 \(u(e_i) = e_{i + 1}\) （\(i\in \mathbf{N}\) ）的自同态。设 \(\mathrm{F} = \mathrm{E}\oplus \mathrm{E}\) ，\(v = u\oplus (u^{*} + 2\cdot 1_{\mathrm{E}})\in \mathcal{L}(\mathrm{F})\) 。设 \(\overline{v} = \pi (v)\) 是 \(v\) 在卡尔金代数 \(\mathcal{Calk}(\mathrm{F})\) 中的像。
\end{enumerate}

\begin{enumerate}
\def\labelenumi{\alph{enumi})}
\item
  自同态 \(v\) 是指标为 -1 的弗雷德霍姆自同态，而 \(v - 2 \cdot 1_{\mathrm{E}}\) 是指标为 1 的弗雷德霍姆自同态。
\item
  设 \(p = \mathrm{X}(\mathrm{X} - 2) \in \mathbf{C}[\mathrm{X}]\) ，S 是 \(\overline{v}\) 的指数谱（I, p. 173 的习题 5）。我们有 \(0 \in p(\mathrm{S})\) ，但 0 不属于 \(p(\overline{v})\) 的指数谱。（利用习题 3。）
\end{enumerate}

